% Part: computability
% Chapter: computability-theory
% Section: equiv-ce-defs

\documentclass[../../../include/open-logic-section]{subfiles}

\begin{document}

\olfileid{cmp}{thy}{eqc}

\olsection[تعريفات المجموعات القابلة للتعداد حسابيًا]{تعريفات متكافئة
  للمجموعات القابلة للتعداد حسابيًا}

لقابلية المجموعة للتعداد حسابيًا أوصاف متكافئة مهمة، تجمعها
المبرهنة الآتية.

\begin{thm}
\ollabel{thm:ce-equiv}
إذا كانت ${\OLId{latin-looped}{S}}$ مجموعة أعداد طبيعية، فالعبارات الآتية متكافئة:
\begin{enumerate}
\item\ollabel{case:ce} ${\OLId{latin-looped}{S}}$ قابلة للتعداد حسابيًا.
\item\ollabel{case:ran-pc} ${\OLId{latin-looped}{S}}$ مدى دالة \emph{جزئية} قابلة للحساب.
\item\ollabel{case:ran-prim} ${\OLId{latin-looped}{S}}$ خالية أو مدى دالة عودية بدائية.
\item\ollabel{case:ce-domain} ${\OLId{latin-looped}{S}}$ \emph{مجال} دالة جزئية قابلة للحساب.
\end{enumerate}
\end{thm}

\begin{explain}
معنى البنود الثلاثة الأولى أن المجموعة غير
الخالية القابلة للتعداد حسابيًا لا يختلف أمرها إذا عددناها بدالة قابلة للحساب، أو جزئية
قابلة للحساب، أو عودية بدائية. وأما البند الرابع فيقرر أنه متى كانت
${\OLId{latin-looped}{S}}$ قابلة للتعداد حسابيًا، فإنه يوجد فهرس~${\OLId{latin-ordinary}{e}}$ بحيث
\[
{\OLId{latin-looped}{S}} = \Setabs{{\OLId{latin-ordinary}{x}}}{\cfind{{\OLId{latin-ordinary}{e}}}({\OLId{latin-ordinary}{x}}) \fdefined}.
\]
فـ${\OLId{latin-looped}{S}}$ تجمع المداخل التي ينتهي عندها حساب
$\cfind{{\OLId{latin-ordinary}{e}}}$. ومن أجل ذلك يقال لهذه المجموعات أحيانًا
\emph{شبه قابلة للقرار}: إن كان العدد منها جاءك بعد حين جواب
«نعم»، وإن لم يكن منها لم يأتك جواب «لا».{}
\end{explain}

\begin{proof}
كل عودية بدائية قابلة للحساب، وكل قابلة للحساب
جزئية قابلة للحساب؛ ولذلك \olref{case:ran-prim} تستلزم
\olref{case:ce}، وتستلزم \olref{case:ce} بدورها
~\olref{case:ran-pc}. (وأما إذا كانت ${\OLId{latin-looped}{S}}$ خالية، فـ${\OLId{latin-looped}{S}}$ مدى الدالة
الجزئية القابلة للحساب غير المعرّفة في أي موضع.) فلا يبقى، إذا أثبتنا أن
\olref{case:ran-pc} تستلزم \olref{case:ran-prim}، شيء لإتمام تكافؤ
البنود الثلاثة الأولى.

فلنفرض أن ${\OLId{latin-looped}{S}}$ مدى الدالة الجزئية القابلة للحساب $\cfind{{\OLId{latin-ordinary}{e}}}$. إن
كانت ${\OLId{latin-looped}{S}}$ خالية فقد تم المطلوب؛ وإلا أخذنا ${\OLId{latin-ordinary}{a}}$ من~${\OLId{latin-looped}{S}}$. وتجيز مبرهنة
كليني للصورة السوية أن نكتب
\[
\cfind{{\OLId{latin-ordinary}{e}}}({\OLId{latin-ordinary}{x}}) = {\OLId{latin-looped}{U}}(\umin{{\OLId{latin-ordinary}{s}}}{{\OLId{latin-looped}{T}}({\OLId{latin-ordinary}{e}}, {\OLId{latin-ordinary}{x}}, {\OLId{latin-ordinary}{s}})}).
\]
وبوجه خاص، تكون $\cfind{{\OLId{latin-ordinary}{e}}}({\OLId{latin-ordinary}{x}})\fdefined$ و$={\OLId{latin-ordinary}{y}}$ إذا وفقط إذا وجد ${\OLId{latin-ordinary}{s}}$
بحيث تصدق ${\OLId{latin-looped}{T}}({\OLId{latin-ordinary}{e}},{\OLId{latin-ordinary}{x}},{\OLId{latin-ordinary}{s}})$ ويكون ${\OLId{latin-looped}{U}}({\OLId{latin-ordinary}{s}})={\OLId{latin-ordinary}{y}}$. فلنعرّف ${\OLId{latin-ordinary}{f}}({\OLId{latin-ordinary}{z}})$ بوضع
\[
{\OLId{latin-ordinary}{f}}({\OLId{latin-ordinary}{z}}) = \begin{cases}
  {\OLId{latin-looped}{U}}(({\OLId{latin-ordinary}{z}})_{\OLMathNumeral{1}}) & \text{إذا صدقت ${\OLId{latin-looped}{T}}({\OLId{latin-ordinary}{e}}, ({\OLId{latin-ordinary}{z}})_{\OLMathNumeral{0}}, ({\OLId{latin-ordinary}{z}})_{\OLMathNumeral{1}})$} \\
  {\OLId{latin-ordinary}{a}}        & \text{في غير ذلك.}
\end{cases}
\]
عندئذ تكون ${\OLId{latin-ordinary}{f}}$ عودية بدائية لأن ${\OLId{latin-looped}{T}}$ و${\OLId{latin-looped}{U}}$ كذلك.
\iftag{TMs}{وبعبارة آلات تورنغ، إذا كان ${\OLId{latin-ordinary}{z}}$ يرمز زوجًا
$\tuple{({\OLId{latin-ordinary}{z}})_{\OLMathNumeral{0}},({\OLId{latin-ordinary}{z}})_{\OLMathNumeral{1}}}$ بحيث تكون $({\OLId{latin-ordinary}{z}})_{\OLMathNumeral{1}}$ حسابًا متوقفًا للآلة~${\OLId{latin-looped}{M}}_{\OLId{latin-ordinary}{e}}$
عند الدخل $({\OLId{latin-ordinary}{z}})_{\OLMathNumeral{0}}$، أعادت ${\OLId{latin-ordinary}{f}}$ خرج الحساب؛ وإلا أعادت~${\OLId{latin-ordinary}{a}}$.}{}

% تصحيح مصدر (C239-01): أُثبت الفرع الثالث الخالي للأمر الثلاثي \iftag.
\textbf{تنبيه تصحيحي على الأصل.} أُثبت الفرع الثالث الخالي للأمر
الشرطي؛ فالأمر ثلاثي الحجج، وإسقاط حجة منه يبتلع ما بعده عند
التوسيع.

والمطلوب أن ${\OLId{latin-looped}{S}}$ عين مدى~${\OLId{latin-ordinary}{f}}$؛ أي إن كل عدد طبيعي~${\OLId{latin-ordinary}{y}}$ يحقق ${\OLId{latin-ordinary}{y}}\in {\OLId{latin-looped}{S}}$
إذا وفقط إذا وقع في مدى~${\OLId{latin-ordinary}{f}}$. أما الوجه الأول، فنفرض ${\OLId{latin-ordinary}{y}}\in {\OLId{latin-looped}{S}}$.
عندئذ يكون ${\OLId{latin-ordinary}{y}}$ في مدى $\cfind{{\OLId{latin-ordinary}{e}}}$، ولذلك توجد قيمتان ${\OLId{latin-ordinary}{x}}$ و~${\OLId{latin-ordinary}{s}}$ تصدق
عندهما ${\OLId{latin-looped}{T}}({\OLId{latin-ordinary}{e}},{\OLId{latin-ordinary}{x}},{\OLId{latin-ordinary}{s}})$ ويكون ${\OLId{latin-looped}{U}}({\OLId{latin-ordinary}{s}})={\OLId{latin-ordinary}{y}}$. لكن عندئذ
${\OLId{latin-ordinary}{y}}={\OLId{latin-ordinary}{f}}(\tuple{{\OLId{latin-ordinary}{x}},{\OLId{latin-ordinary}{s}}})$. وأما العكس، فلنأخذ ${\OLId{latin-ordinary}{y}}$ في مدى~${\OLId{latin-ordinary}{f}}$. عندئذ إما
${\OLId{latin-ordinary}{y}}={\OLId{latin-ordinary}{a}}$، وإما أنه يوجد~${\OLId{latin-ordinary}{z}}$ بحيث تصدق ${\OLId{latin-looped}{T}}({\OLId{latin-ordinary}{e}},({\OLId{latin-ordinary}{z}})_{\OLMathNumeral{0}},({\OLId{latin-ordinary}{z}})_{\OLMathNumeral{1}})$ ويكون
${\OLId{latin-looped}{U}}(({\OLId{latin-ordinary}{z}})_{\OLMathNumeral{1}})={\OLId{latin-ordinary}{y}}$. ولما كانت، في الحالة الأخيرة،
$\cfind{{\OLId{latin-ordinary}{e}}}(({\OLId{latin-ordinary}{z}})_{\OLMathNumeral{0}})\fdefined={\OLId{latin-ordinary}{y}}$، فإن ${\OLId{latin-ordinary}{y}}$ في~${\OLId{latin-looped}{S}}$ في الحالتين.

(يعني الترميز $\cfind{{\OLId{latin-ordinary}{e}}}({\OLId{latin-ordinary}{x}})\fdefined={\OLId{latin-ordinary}{y}}$ أن «$\cfind{{\OLId{latin-ordinary}{e}}}({\OLId{latin-ordinary}{x}})$ معرفة
وتساوي ${\OLId{latin-ordinary}{y}}$». ولنا أن نكتفي بكتابة $\cfind{{\OLId{latin-ordinary}{e}}}({\OLId{latin-ordinary}{x}})={\OLId{latin-ordinary}{y}}$ بالقدر نفسه، غير أن
زيادة السهم تنبّه أحيانًا إلى جزئية الدالة.)

وبقي من برهان \olref{thm:ce-equiv} إثبات تكافؤ
\olref{case:ce} و~\olref{case:ce-domain}. نبدأ بإثبات أن
\olref{case:ce} تستلزم~\olref{case:ce-domain}. أما المجموعة الخالية فهي مجال الدالة غير المعرّفة في أي موضع. وأما غير الخالية فنفرض ${\OLId{latin-looped}{S}}$ مدى دالة
قابلة للحساب~${\OLId{latin-ordinary}{f}}$، أي
\[
{\OLId{latin-looped}{S}} = \Setabs{{\OLId{latin-ordinary}{y}}}{\text{لبعض ${\OLId{latin-ordinary}{x}}$، } {\OLId{latin-ordinary}{f}}({\OLId{latin-ordinary}{x}}) = {\OLId{latin-ordinary}{y}}}.
\]
ولتكن
\[
{\OLId{latin-ordinary}{g}}({\OLId{latin-ordinary}{y}}) = \umin{{\OLId{latin-ordinary}{x}}}{({\OLId{latin-ordinary}{f}}({\OLId{latin-ordinary}{x}}) = {\OLId{latin-ordinary}{y}})}.
\]
فتكون ${\OLId{latin-ordinary}{g}}$ جزئية قابلة للحساب، ولا تتعين ${\OLId{latin-ordinary}{g}}({\OLId{latin-ordinary}{y}})$ إلا حين
يصدق ${\OLId{latin-ordinary}{f}}({\OLId{latin-ordinary}{x}})={\OLId{latin-ordinary}{y}}$ لعدد ما~${\OLId{latin-ordinary}{x}}$، وتتعين كلما صدق. فمجال~${\OLId{latin-ordinary}{g}}$ هو بعينه مدى~${\OLId{latin-ordinary}{f}}$.
\iftag{TMs}{وبعبارة آلات تورنغ: إذا أعطيت آلة تورنغ~${\OLId{latin-looped}{F}}$ تعدد عناصر
~${\OLId{latin-looped}{S}}$، فلتكن ${\OLId{latin-looped}{G}}$ آلة تورنغ شبه تقرر ${\OLId{latin-looped}{S}}$ بالبحث في مخارج~${\OLId{latin-looped}{F}}$ لترى هل
عنصر معطى في المجموعة؛ فتتوقف إذا كان فيها، وتواصل البحث إلى الأبد إذا
لم يكن فيها.}{}

وأما كون \olref{case:ce-domain} مستلزمًا لـ~\olref{case:ce}، فبيانه أن نفرض
أن ${\OLId{latin-looped}{S}}$ مجال الدالة الجزئية القابلة للحساب~$\cfind{{\OLId{latin-ordinary}{e}}}$، أي
\[
{\OLId{latin-looped}{S}} = \Setabs{{\OLId{latin-ordinary}{x}}}{\cfind{{\OLId{latin-ordinary}{e}}}({\OLId{latin-ordinary}{x}}) \fdefined}.
\]
فإن كانت ${\OLId{latin-looped}{S}}$ خالية تم المطلوب؛ وإلا أخذنا ${\OLId{latin-ordinary}{a}}$ من~${\OLId{latin-looped}{S}}$، وعرّفنا ${\OLId{latin-ordinary}{f}}$ بوضع
\[
{\OLId{latin-ordinary}{f}}({\OLId{latin-ordinary}{z}}) = \begin{cases}
({\OLId{latin-ordinary}{z}})_{\OLMathNumeral{0}} & \text{إذا صدقت ${\OLId{latin-looped}{T}}({\OLId{latin-ordinary}{e}},({\OLId{latin-ordinary}{z}})_{\OLMathNumeral{0}},({\OLId{latin-ordinary}{z}})_{\OLMathNumeral{1}})$} \\
{\OLId{latin-ordinary}{a}} & \text{في غير ذلك.}
\end{cases}
\]
وعندئذ، كما سبق، يكون العدد ${\OLId{latin-ordinary}{x}}$ في مدى~${\OLId{latin-ordinary}{f}}$ إذا وفقط إذا كانت
$\cfind{{\OLId{latin-ordinary}{e}}}({\OLId{latin-ordinary}{x}})\fdefined$، أي إذا وفقط إذا كان ${\OLId{latin-ordinary}{x}}\in {\OLId{latin-looped}{S}}$.
\iftag{TMs}{وبعبارة آلات تورنغ: إذا أعطيت آلة ${\OLId{latin-looped}{M}}_{\OLId{latin-ordinary}{e}}$ شبه تقرر ${\OLId{latin-looped}{S}}$،
فعدد عناصر ${\OLId{latin-looped}{S}}$ بالمرور على جميع حسابات آلات تورنغ الممكنة وإعادة
المداخل التي تقابل حسابات متوقفة.}{}
\end{proof}

ومن البند~\olref{case:ce-domain} في \olref{thm:ce-equiv} نحصل على طريقة
لتعداد المجموعات القابلة للتعداد حسابيًا: نجعل، لكل~${\OLId{latin-ordinary}{e}}$، الرمز ${\OLId{latin-looped}{W}}_{\OLId{latin-ordinary}{e}}$ لمجال
$\cfind{{\OLId{latin-ordinary}{e}}}$، أي
\[
{\OLId{latin-looped}{W}}_{\OLId{latin-ordinary}{e}} = \Setabs{{\OLId{latin-ordinary}{x}}}{\cfind{{\OLId{latin-ordinary}{e}}}({\OLId{latin-ordinary}{x}}) \fdefined}.
\]
وعندئذ، إذا كانت ${\OLId{latin-looped}{A}}$ أي مجموعة قابلة للتعداد حسابيًا، كان
${\OLId{latin-looped}{A}}={\OLId{latin-looped}{W}}_{\OLId{latin-ordinary}{e}}$ لبعض~${\OLId{latin-ordinary}{e}}$.

وللمجموعات القابلة للتعداد حسابيًا توصيف آخر كما يأتي.

\begin{thm}
\ollabel{thm:exists-char}
تكون المجموعة ${\OLId{latin-looped}{S}}$ قابلة للتعداد حسابيًا إذا وفقط إذا وجدت علاقة قابلة
للحساب ${\OLId{latin-looped}{R}}({\OLId{latin-ordinary}{x}},{\OLId{latin-ordinary}{y}})$ بحيث
\[
{\OLId{latin-looped}{S}} = \Setabs{ {\OLId{latin-ordinary}{x}} }{ \lexists[{\OLId{latin-ordinary}{y}}][{\OLId{latin-looped}{R}}({\OLId{latin-ordinary}{x}},{\OLId{latin-ordinary}{y}})] }.
\]
\end{thm}

\begin{proof}
أما الاتجاه الأول، فلنفرض قابلية ${\OLId{latin-looped}{S}}$ للتعداد حسابيًا. فيكون
${\OLId{latin-looped}{S}}={\OLId{latin-looped}{W}}_{\OLId{latin-ordinary}{e}}$ لبعض ${\OLId{latin-ordinary}{e}}$. ولهذه القيمة من~${\OLId{latin-ordinary}{e}}$ يمكننا كتابة ${\OLId{latin-looped}{S}}$ هكذا:
\[
{\OLId{latin-looped}{S}} = \Setabs{ {\OLId{latin-ordinary}{x}} }{ \lexists[{\OLId{latin-ordinary}{y}}][{\OLId{latin-looped}{T}}({\OLId{latin-ordinary}{e}}, {\OLId{latin-ordinary}{x}}, {\OLId{latin-ordinary}{y}})] }.
\]
وأما الاتجاه العكسي، فنفرض
${\OLId{latin-looped}{S}}=\Setabs{{\OLId{latin-ordinary}{x}}}{\lexists[{\OLId{latin-ordinary}{y}}][{\OLId{latin-looped}{R}}({\OLId{latin-ordinary}{x}},{\OLId{latin-ordinary}{y}})]}$. عرف ${\OLId{latin-ordinary}{f}}$ بـ
\[
{\OLId{latin-ordinary}{f}}({\OLId{latin-ordinary}{x}}) \simeq \umin{{\OLId{latin-ordinary}{y}}}{{\OLId{latin-looped}{R}}({\OLId{latin-ordinary}{x}}, {\OLId{latin-ordinary}{y}})}.
\]
عندئذ تكون ${\OLId{latin-ordinary}{f}}$ جزئية قابلة للحساب، وتكون ${\OLId{latin-looped}{S}}$ مجال~${\OLId{latin-ordinary}{f}}$.
\end{proof}


\end{document}
